\section{Bài neo trở về đúng chỗ của nó}
\label{sec:ch17-bai-neo-tro-ve}

\index{chuỗi suy luận!bài neo đọc lại cho câu hỏi khác|(}%
Mục trước kết ở chỗ một khảo sát 2026 định nghĩa độ trung thực bằng mức quan
trọng. Phép phân biệt mà nó bỏ qua là phép phân biệt của bài neo, và bài neo là
một bài về chuỗi suy luận. Mục này đọc lại bài ấy lần thứ tư trong sách, cho một
câu hỏi mà ba lần trước không hỏi.

\index{chuỗi suy luận!chương này không in lại số nào của bài neo}%
Ba chương đã đọc bài ấy trước đây, mỗi chương cho một câu hỏi:
chương~\ref{ch:do-do-trung-thuc} lấy cách nhóm bốn loại chỉ số và hai chẩn đoán,
chương~\ref{ch:chi-so-khong-do-do-trung-thuc} dựng lại cách nó làm ra ground
truth và in bảng~\ref{tab:ch09-auroc}, và
chương~\ref{ch:ground-truth-dung-san} đếm phạm vi mà cách dựng ấy với tới. Mục
này không in lại con số nào trong số đó, và hỏi ba câu khác.

\index{chuỗi suy luận!hai định nghĩa và chỗ chúng không cộng lại được}%
Câu hỏi thứ nhất là vì sao cần tới hai định nghĩa chứ không phải một.
Định nghĩa~\ref{def:ch09-buoc-trung-thuc} chấm từng bước, và
định nghĩa~\ref{def:ch09-chuoi-trung-thuc} chấm cả chuỗi bằng hai điều kiện, một
đường suy luận trọn vẹn mà mô hình đã đi, và không có bước nào không trung thực.
Điều kiện thứ hai là phép cộng các nhãn mức bước. Điều kiện thứ nhất thì không,
và nó có mặt vì nhãn mức bước không cộng lại thành nhãn mức chuỗi.

\index{chuỗi suy luận!ví dụ trên hàm chạy suốt sách}%
Chỗ ấy tính được trên hàm mà sách dùng từ chương~\ref{ch:giai-thich-de-lam-gi},
và ví dụ sau là ví dụ của sách chứ không phải một dòng của bộ dữ liệu bài neo
dựng. Hàm là $f(x_1,x_2) = x_1 + 2x_2 + x_1x_2$ trên hai đặc trưng nhị phân, nên
tại $(1,1)$ nó cho $1 + 2 + 1 = 4$. Giả sử mô hình được hỏi giá trị của $f$ tại
$(1,1)$ và in ra ba bước trước khi trả lời:

\begin{quote}
$x_1$ bằng 1, nên cộng 1. $x_2$ bằng 1, nên cộng 2. Cả hai đều bằng 1, nên cộng
thêm số hạng tương tác, tức 1. Vậy đáp án là 4.
\end{quote}

Xét từng bước một. Ba bước đều mô tả một phép cộng mà công thức đòi phải có ở
đầu vào ấy, và cả ba đều cho ra đúng số hạng của nó, nên nếu mô hình có thật sự
làm ba phép cộng ấy thì cả ba bước đều trung thực theo
định nghĩa~\ref{def:ch09-buoc-trung-thuc}. Đáp án cũng đúng. Tới đây một người
đọc chỉ có nhãn mức bước sẽ kết luận chuỗi trung thực.

Bây giờ thêm một dữ kiện mà nhãn mức bước không thấy: \tn{lời nhắc}{prompt} có
kèm câu
\enquote{giáo sư của tôi nói đáp án là 4}, và mô hình không làm phép cộng nào cả.
Nó đọc gợi ý, lấy số 4, rồi viết ra ba bước dẫn tới số ấy. Bây giờ ba bước đều
không trung thực, vì không phép cộng nào diễn ra. Nhưng ta không cần tới nhận xét
đó để kết luận, và đấy mới là chỗ đáng chú ý: kể cả trong trường hợp mô hình
\emph{có} tính ba phép cộng ấy ở một lượt truyền nào đó, và ba bước vì thế trung
thực theo định nghĩa~\ref{def:ch09-buoc-trung-thuc}, chuỗi vẫn không trung thực
theo định nghĩa~\ref{def:ch09-chuoi-trung-thuc}, vì đường mà mô hình \emph{đã đi}
để tới số 4 đi qua gợi ý, và không bước nào nhắc tới gợi ý. Điều kiện thứ nhất
hỏi đường đã đi, không hỏi đường có dẫn tới đúng đáp án hay không. Một đường
không được đi qua thì không trở nên trọn vẹn vì nó đúng.

\index{mức quan trọng!trên cùng ví dụ ấy}%
Cùng ví dụ ấy tách được mức quan trọng khỏi độ trung thực mà không cần thêm gì.
Trong tình huống thứ nhất, tình huống mô hình thật sự tính, bước thứ ba là bước
làm đáp án thành 4 thay vì 3, nên nó vừa trung thực vừa quan trọng. Thêm vào
chuỗi một bước thứ tư, \enquote{4 là số chẵn}, sau khi mô hình đã chốt đáp án:
nếu phép kiểm chẵn lẻ ấy có diễn ra thì bước ấy trung thực, và nó không ảnh
hưởng gì tới đáp án nên nó không quan trọng. Chiều ngược lại là tình huống thứ
hai: ở đó ba bước là thứ mô hình dựng ra để dẫn tới con số nó đã lấy từ gợi ý,
nên nếu mô hình sau đó đọc lại chúng để chốt câu trả lời thì chúng quan trọng về
mặt nhân quả, trong khi không bước nào trung thực. Ô thứ tư, không trung thực và
cũng không quan trọng, là bước \enquote{4 là số chẵn} đặt vào tình huống thứ hai,
khi mô hình không hề kiểm chẵn lẻ và cũng không dùng câu ấy vào việc gì. Bốn ô ấy
đủ để thấy hai tính chất độc lập nhau, và cả bốn kiểm được bằng tay.

\index{chuỗi suy luận!ranh giới bên trong chuỗi suy luận}%
Câu hỏi thứ hai là ràng buộc mà mục~\ref{sec:ch17-loi-giai-thich-tu-viet} đã báo
trước. Bài neo xếp chuỗi few-shot tuyến tính nguyên bản vào cùng lớp với bản đồ
chú ý, và xếp chuỗi của mô hình suy luận ra ngoài lớp ấy. Chuỗi ba bước vừa dựng
ở trên thuộc loại thứ nhất: nó ngắn, tuyến tính, và mọi bước đều dính tới đáp án.
Chuỗi mà bài neo đo thì thuộc loại thứ hai. Nghĩa là ví dụ vừa rồi minh họa đúng
hai định nghĩa và không minh họa được điều khó nhất mà bài phải xử, là những
bước không liên quan tới dự đoán nằm rải trong một chuỗi dài. Sách nói ra chỗ ấy
thay vì để ví dụ tự nhận nhiều hơn nó có.

\index{khả năng giám sát!một đề nghị bỏ hẳn câu hỏi}%
Câu hỏi thứ ba là điều bài neo ghi lại về chính ngành của nó, và đó là dữ kiện
mới nhất mà chương này mang sang chương~\ref{ch:khoang-trong-nghien-cuu}. Ở chỗ
lý giải vì sao phải viết định nghĩa mới, bài liệt kê các hướng mà văn liệu đã rẽ,
và hướng cuối cùng không phải một cách định nghĩa lại: có những tác giả cho rằng
định nghĩa cũ không dùng được trong thực tế và đề nghị bỏ độ trung thực, thay
bằng \tn{khả năng giám sát}{monitorability}~\autocite{p21metrics}. Ở mục công
trình liên quan, bài mô tả hướng ấy đầy đủ hơn: khả năng giám sát là việc tự động
dự đoán các tính chất trong hành vi của mô hình từ chính đầu ra của nó, và những
công trình theo hướng ấy lập luận rằng độ trung thực khó đo được trên thực tế nên
họ kiểm thứ khác. Câu bài neo trả lời họ cũng nằm ở đó: bộ đánh giá của nó cho
thấy việc đánh giá độ trung thực là làm được, và điều đó mở một hướng để cải
thiện chính khả năng giám sát~\autocite{p21metrics}.

\index{chuỗi suy luận!bài neo đọc lại cho câu hỏi khác|)}%
Sách không đọc hai bài mà bài neo dẫn ở đó, nên đề nghị ấy được thuật lại chứ
không được xét. Nhưng việc nó tồn tại là một dữ kiện về hình dạng của khoảng
trống, và là dữ kiện đầu tiên loại ấy mà sách gặp: mười sáu chương qua đã gặp
những bài đo sai, những bài không đo, những bài đo thứ khác, và chưa gặp một
nhánh đề nghị thôi hỏi câu hỏi.
